\documentclass{article}
\usepackage{amsmath}
\newcommand{\argmin}{\arg\!\min} 
\newcommand{\argmax}{\arg\!\max} 
\usepackage{amssymb}
\usepackage{amsfonts}
\usepackage[T1]{fontenc}
\usepackage{bm}
\usepackage{array}
\usepackage{graphicx}
\usepackage[utf8]{inputenc}

\title{Drawing demo}
\date{\today}
\author{Marko Guberina}

\begin{document}
\maketitle

\section{Introduction}
The demo is for the robot to replicate a drawing you make on the computer.
The complete demo works as follows:
\begin{enumerate}
		\item You are prompted with a blank canvas on which you are to draw an image.
				You are presumably close to the robot.
		\item The robot then comes to you to pick up a marker with which to draw.
		\item The robot calibrates the length of the marker (the grasp and marker size are unknown)
				by touching the board once.
		\item The robot then proceeds to create the drawing. It treats the drawing trajectory
				as a dynamic motion primitive (DMP). 
				It uses impedance to maintain contact with the whiteboard.
		\item The robot comes back to pick up an eraser.
		\item It calibrates the length of the eraser.
		\item It then "draws" a wiping picture. But the stiffness is higher, 
				and the target is now more into the board.
\end{enumerate}

The main file is drawing\_from\_input\_drawing.py.

\section{How it works on 2025-11-26}
\subsection{Things to be done before you can run the demo}
\subsubsection{Find a whiteboard and a few markers}
Get a whiteboard and some whiteboard markers. Preferably an old whiteboard out of use,
as it is very possible (likely even) the robot will make dents in it.
Also, the robot will certainly ruin 1 marker, and there's a high possibility it will be more than one.
This is because, at least on UR5e, the force-torque sensor is not very sensitive,
so we have to shove the marker into the board with forces between 2.5-5N (the force isn't controlled directly,
but more on this later). 
And the drawing is done with the marker being perpendicular.

The reasons behind this is that we don't assume great board calibration nor that the board is perfectly flat.
At least ours isn't.

\subsubsection{Manually define a sensible work pose for the robot}
Basically find a home position from which Cartesian controllers can easily
access the board. This is very important because the demo does not
account for any kind of collision constraints, it just assumes
that everything will be fine. From experience there were no problems
with this as the robot moves relatively little to do the drawing.

Write this joint configuration down.

\subsubsection{Marker pickup/handover}
To add some flare to the demo, 
there is a feature for the robot to come to you so that you give it the marker,
or cloth to wipe the board with.
The idea being that you sit behind a computer close to the robot. It would be
very easy to do a pick and place instead (feel free to put this in),
this is just what's in the code right now.

The handover is executed as joint-space motion, so you need to provide 
joint angles constituting the path.
The current expected way to do this is to 
run the convenience\_toolbox/freedrive\_v2.0.py .
It puts the robot to freedrive/leadthrough and allows you to save 
angles by pressing a button on the keyboard.
The initial pose for the robot should be the "drawing home" you defined in the previous step.
These angles are then fed into a DMP which executes the trajectory.
The gripper then opens, waits a few seconds, and then closes.
The same exact trajectory, only reversed, is executed to go back.
NOTE that you need to manually rename the file outputed from freedrive
to ./data/from\_writing\_to\_handover.pickle\_save. Yes, 
this should be in parameters, and yes this step should be automated
(or at least called from the main file).

\paragraph{Note} This option is off by default. Enable it with --pick-up-marker.

\subsubsection{Whiteboard (plane) pose calibration}
Before any drawing, board calibration needs to be performed. Do this with --calibrate-plane
on the main file.
The program will guide to the setup, but the point is you go into freedrive
and put the end-effector to the top-left corner of the board.
The program than samples from the board point to the right and down
by touch the board. It updates the orientation estimate of the end-effector as it goes,
so that you (eventually) get a flat contact point to the board.
It then applies a linearly weighted least-squares fit to the accumulated points,
biasing the more recent samples.

Of course, you need to run calibration only once. The pose is saved
into ./parameters/plane\_pose.pickle. 
At the moment the actual file being used is ./parameters/plane\_pose.pickle\_save,
meaning you have to manually rename it. This is a stupid atavism and will 
be removed asap (no promises at the time of writing though).

\subsubsection{Do a few runs and tune the parameters}
The defaults are sane, but the particulars always depend
on the concrete robot and other equipment used.
What to pay attention to is described in the following section
as it was more natural to write that way.


\subsection{Normal demo operation}
\textbf{First you are prompted to draw the image}. It has to be a single
continuous line. You can redo the drawing by just attempting to draw
a new one. Accept by pressing Enter.
You can also reuse the previous drawing if you selected --no-draw-new
(this is handy for board wiping, more on that later).

\textbf{The next action is the optional handover/pickup}.
It then possibly first does a moveJ to home, not sure.
\textbf{The robot then goes to a point above  the board} (as in away from the board) which was the first 
calibration point. \textbf{It then go straight down into the board
to touch it once}, and it immediately goes back up to the same point.
This is done to calibrate the length of the marker which is unknown
both due to the handover, and, even if you didn't go for the handover,
due to the possible in-hand slippage during writing (it can and does happen,
but not super often - the issue is that this number need to be precise).

With this, we can map the drawing onto the board, i.e. create the SE3 reference
out of it.
Once this is done, a simulation is run to moveL to every point of the reference
and record the joint angles. This takes about 1-10 seconds depending on path length.
\textbf{IMPORTANT TO UNDERSTAND POSSIBLE FAILURES} From these joint angles a DMP is formed. 
The timing for the DMP is provided with --tau0.
There is automatic scaling if the robot can't make it in time.
But this distorts the image quite a bit. By this I mean it does not work at all,
but at least the robot does not explode (as in do something random and violent). 
Anyhow, a 10 second default is chosen.
This is enough for basic shapes which take little time to draw,
which is what is assumed to be default (I personally always draw hearts).
Other possible distortions are due to board bendiness and bad f/t readings.

\textbf{We then do a compliantMoveL to the beginning of the drawing}.
\textbf{Drawing starts immediately after}. Impedance
in the form of direct force feedback is added to the DMP.
Only the end-effector z axis is compliant, the rest is fully stiff.
The reference path is inside the board. How much is determined
with --mm-into-board, with 3mm being the default.
In our experiments this results in 2.5-5N of force.
Look at the wrench plot to see how much you use (NOTE that the wrench\_used i.e. filtered looks bigger
as we multiply the rotational dofs as 1-10 linear-rotational ratio is
generally nice compliance). Remember to use
 --save-log, --index-runs and --run-name to save the log as they are not saved by default.
 This way you can analyse this in peace, instead of while the robot is running.
\textbf{IMPORTANT TO UNDERSTAND POSSIBLE FAILURES} If the robot does something
weird like randomly push to hard into the board, and then go back,
or just go back, this is almost certainly a faulty f/t sensor reading.
While you can make this better with better filtering, be wary
of garbage in, garbage out - this is the only sensing we use in the demo,
and if this is off there's fundamentally no way to save it.

\textbf{Finally, the robot does compliantMoveL to go above the board.
The program then exits.}

However, this is not necessarily the end of the demo!
\textbf{You can re-run the script, but with
--pick-up-marker --board-wiping --mm-into-board=XY --no-draw-new 
to pick up a cloth to wipe the board with, and wipe the drawing}
by tracing the exact drawing path back! Can't be more efficient than that.
How much --mm-into-board you need depends on the cloth used.
We got the best results with a strongly folded set of kitchen towels
as this was just compliant enough (so it's rather sturdy, but not a rigid body).
You might also experiment with --alpha which determines
the amount of compliance (higher means more compliance),
but in our experience just changing the reference was enough.


\section{Limitations}
\begin{itemize}
		\item The drawing has to be a single continuous line. 
				\begin{itemize}
						\item A different program for drawing is necessary (or the current one needs to be extended).
						\item The same controller could be applied for each line, but logic to go line by line
								needs to be implemented. 
						\item To avoid it being excruciatingly slow, the default timing for a line needs
								to be replaced with a heuristic.
				\end{itemize}
		\item The change of the reference due to impedance distorts the picture. Intricate features
				won't happen. A different controller would need to be written to more aggressively favour
				path following over timing (or disregard timing entirely).
		\item The performance depends on hardware and on calibration a bit more than it should.
				Some of the calibration is automated, but one still needs to devote around 30min
				to tune all the parameters for the best performance.
\end{itemize}


Things to improve/re-solve the problem:
\begin{enumerate}
		\item Allow the viz option to draw every path point.
		\item Add a non-convergence detection to compliantMoveL, so that if you didn't move much,
				and you're already close to the goal, just exit - you will never converge 
				due to the f/t noise being stronger than the signal to converge.
				So \fbox{\texttt{if err < 0.01 and avg(v) < 0.01: return}}
		\item Switch DMP to Cartesian space. This way you don't need to 
				recompute everything to joint space - faster!
		        BONUS: you can skip the offset, and just get it from the first
				point to be drawn --- then update the rest of the reference.
		\item Implement better kinesthetic teaching so that the handover is easier to do.
		\item Try doing it with basic Cartesian path following instead (skip the DMP).
		\item Optimize moveLs so that they're faster (the default point-to-point
				velocity profile is kinda ass).
\end{enumerate}

Micro quality of life improvements (maybe):
\begin{itemize}
		\item Replace tau0 (DMP traj) from an argument into a heuristic depending on the length
				of the trajectory. Anything sensible will be better than the current version
				which already is a heuristic and only works for short drawings.
		\item Add an automatic restart if you observe a f/t sensor crap out.
\item Handover joint angles definition via kinesthetic teaching should be added into the main program with an --initialize-demo argument
which is False by default. It would also automatically set the calibration flag to True.
\item In retrospect, it would be better if the calibration process
starts with the end-effector touching the board, and then go up a bit,
instead of starting from the top. This way the initial estimate is better.
Also, maybe it would be better to simply discard the initial points, idk.
It does work so this is kind of pointless optimization.
\item The output of the 
\end{itemize}

\end{document}
